\documentclass[11pt, a4paper, UTF8]{chinese-erj}

% --- 必须的宏包 ---
\usepackage{amsmath, amssymb, amsthm, mathtools}
\usepackage{booktabs} % 用于绘制三线表
\usepackage{float}    % 用于控制表格位置
\usepackage{url}      % 用于处理超链接与长网址自动换行
\usepackage{caption}  % 用于自定义图表标题格式

% --- 强制将 Table 改为 表，并把冒号改为标准空格 ---
\renewcommand{\tablename}{表}
\renewcommand{\figurename}{图}
\captionsetup[table]{labelsep=quad} % 效果：表1 变爻个数的分布

% --- 定义定理环境 ---
\newtheorem{theorem}{定理}
\newtheorem{lemma}{引理}
\renewcommand{\proofname}{\textbf{证明}}
\renewcommand{\qedsymbol}{$\blacksquare$} % 霸气实心黑方块

% --- 调整页眉高度以防警告 ---
\setlength{\headheight}{16pt}
\addtolength{\topmargin}{-4pt}

% --- 动态引入本篇稿件的参考文献库 ---
\addbibresource{AERub_cyberdiv.bib}

% --- 动态定义出版信息 ---
\header{晨韶学妹：赛博算卦为何很准？}
\pubinfo{2026年第2卷第1期} 

\begin{document}

\setcounter{page}{11} 

% ==========================================
% 标题与作者设置 (极简风，合并底层声明)
% ==========================================
\title{赛博算卦为何很准？——基于chgdiv6命令的实证分析\thanks{晨韶学妹，Hsiakuanying University。本文受自燃科学鸡精委重点项目（项目批准号：72636007）“人类迷惑行为的经济学解释”资助。作者感谢匿名审稿专家的宝贵建议，当然文责自负。}}

\author{晨韶学妹}
\date{}

\maketitle
% \thispagestyle{empty}

\begin{abstract}
本文基于chgdiv6命令生成的占卜结果，实证分析“赛博算卦为何很准”的现象。研究发现，求卦者在选择随机数种子时无意识地融入自身心理状态与生活背景：对幸运数字“6”的偏好显著降低了占卜结果的吉利程度；义务教育制度实施后，出生日期决定的入学真实年龄与占卜结果呈显著正相关。研究表明，占卜的“准确性”源于人类行为的可统计规律性，而非超自然力量。

\keywords{赛博算卦 \quad 幸运数字偏好 \quad 出生日期 \quad 相对年龄效应 \quad 有序离散模型}
\end{abstract}

% --- 正文 ---
\section{引言}

2023年9月，兰州大学的陈志豪同学研发出了一个基于Stata的占卜命令chgdiv6，于2025年2月在Gitee与Github上发布\footnote{Gitee 仓库网址：\url{https://gitee.com/chenshaoxuemei/chgdiv6_bkp}}，该命令以有趣的功能在B站\footnote{【Stata也能算卦】chgdiv6命令：在Stata中算铜钱卦 \url{https://www.bilibili.com/video/BV1bGQTYFEte/?share_source=copy_web&vd_source=029ff775c4cddfa96afc77dede9af232}}、小红书、知乎等多个平台上获得了较多的关注。一个有趣的发现是，尽管以设计者的身份看来，这个命令只是用中国传统文化对64位梅森旋转法予以包装，但大多数寻求占卜的人都认为占卜的结果符合他们的预期。本文尝试以简单的有序离散模型分析对该现象予以解释。

chgdiv6命令具备三种起卦方式：（1）以当前时间起卦；（2）指定随机数起卦；（3）输入文字起卦。从本质上而言，以当前时间起卦与输入文字起卦均是基于指定随机数起卦的原理，前者使用组合后的时间为随机数种子，后者使用每个字符转换得到的UTF-8字节序列之和为随机数种子。由于该命令者曾指出，其长期使用指定随机数起卦时观察到了前文所述的现象，故本文主要研究直接指定随机数起卦方式下，占卜为何符合求卦者的预期。

求卦者在选定随机数种子时，存在着以下两种偏好：（1）倾向于选择与自身年龄、生日或重要日期相关的数字；（2）在特定心理预期下选择“幸运数字”。前者在占卜时主要出于兴趣，大多数情况下所问内容是具有长期性质的“运势”。后者在一定心理影响下，占卜所问内容涉及具体的事件，实则反映了特定的心理预期。

关于出生日期如何影响人生发展，已有研究揭示了若干值得关注的现象。在教育与发展心理学领域，这一影响主要通过“相对年龄效应”体现。由于“入学年龄截止日期”的制度安排，同一学年入学的儿童在入学时的实际年龄存在微小差异，这一差异会逐渐累积，在学业表现、社交能力与自信心等方面形成持续影响，甚至可能对个体长期的教育成就和职业选择产生作用\parencite{xu2021}。由此可见，基于出生日期形成的入学真实年龄差异，可以在一定程度上反映个人的长期发展轨迹。由此可见，基于出生日期计算的入学真实年龄，一定程度上能够反映个人的长期运势。

在探讨个体在不同心理状态下对幸运数字的选择偏好方面，大量研究发现，负性情绪显著增强了个体对带有文化吉祥寓意数字的偏好，这被视为一种心理调适机制。在焦虑或无助情境下，数字“8”不仅是一个数学符号，更被赋予了“发”“顺”“圆满”等积极文化内涵，成为一种心理锚点。类似地，数字“6”因其谐音“顺”（顺利、顺畅）而与传统祝福语“六六大顺”相关联，承载着积极的吉祥寓意。例如，\textcite{chen2022} 研究发现，负面情绪可提高消费者对以“8”结尾的“幸运价格”产品的购买意愿；\textcite{luo2022} 则发现，含数字“8”的房间定价不仅能提升预订量，也有助于提高客户评分。然而，当前关于幸运数字的研究仍以消费行为实验为主，故对于幸运数字的应用探讨多与消费决策密切相关，尚缺乏更多元情境的拓展。

关于占卜结果的解读，自孔子编撰《十翼》以来，学界已有大量探讨。《易传》开篇有云：“圣人设卦观象，系辞焉而明吉凶，刚柔相推而生变化。”此句表明，卦爻符号并非静止不变，其所象征的情境亦随之流转，因而人的应对策略也应动态调整，因事制宜。正因如此，《周易》占卜所生成的4096种结果，长期以来并未形成统一的吉凶分类标准。虽有民间学者（如托名邵雍）尝试将64卦分为上上、中上、中中、中下、下下五等，亦或对384爻进行吉凶划分，但这类划分方式始终存在争议。例如，鼎卦卦辞为“元吉，亨”，却被归类为中下卦；而睽卦卦辞“小事吉”则被普遍解读为不吉之象。综上，由于《周易》的解读历来强调因时因地因事制宜，学界既未对其展开明确的吉凶划分，亦缺乏系统量化的研究。为便于操作，本文主要参考民间对384爻及64卦的吉凶划分方式，并采用朱熹提出的占卜方法，对4096种占卜结果进行量化评分。

为探讨基于六位出生日期代码是否能够反映个人长期运势，本文以近100年间共36,525条六位出生日期代码为样本起卦，并依据1986年《中华人民共和国义务教育法》确立义务教育制度的时间节点，将样本分为前后两组，重点考察入学时实际年龄对占卜结果的影响。研究发现，在我国正式实施义务教育后，入学时实际年龄与占卜结果的吉利程度呈显著正相关，且在控制变爻个数后该关系依然成立。在心理状态与幸运数字偏好的维度上，本文发现以数字“6”结尾的随机数种子起卦，会显著降低占卜结果的吉利程度。这一结果与负性情绪增强个体调用文化幸运数字以寻求心理缓冲的既有发现具有一致性。

本文的边际贡献主要体现在以下三个方面：
第一，拓展了数字偏好研究的应用场景。以往对幸运数字的研究多集中于消费行为领域，如价格尾数、房间编号等。本文首次将幸运数字偏好引入数字占卜这一新兴文化现象中，揭示了其在非消费情境下的心理作用机制。
第二，为出生日期的长期影响提供了新的实证证据。本文以占卜结果为因变量，验证了出生日期对个体发展的影响，尤其是在义务教育制度实施后的显著正向关系，为教育心理学中的“相对年龄效应”提供了跨学科的量化支持。
第三，推动了数字文化研究与计量方法的交叉融合。本文采用有序离散模型、固定效应估计与聚类稳健标准误等方法，对传统文化现象进行了系统性的计量分析，展示了社会科学方法在人文学科研究中的适用性与创新性。

\section{数据与实证策略}

本文首先使用 1 至 4,194,304 为随机数种子起卦，得到了 4,194,304 个占卜结果样本。在研究出生日期的影响部分，则利用 1925年1月1日（250101）至 2024年12月31日（261231）的六位出生日期数字为随机数种子起卦，得到 36,525 条占卜结果。

起卦方式为铜钱卦，具体方法为连续六次抛掷三枚硬币，自下而上地依次画出每次结果所对应的爻象，从而构成一个完整的六爻卦象。每次抛掷，根据三枚铜钱正反面的组合，可以产生四种基本爻象：一个背面为少阳，两个背面为少阴，三个背面为老阳，三个正面为老阴。其中，老阳和老阴因其“物极必反”的特性，会引发卦象的变化，即在本卦（主卦）的基础上，将老阳爻变为少阴、老阴爻变为少阳，从而衍生出一个新的卦象，称为变卦（之卦）。不难计算得到，一共有 $2^6$ 即 4096 种结果。同时，不难发现四种基本爻象的概率分别为 $3/8$、$3/8$、$1/8$ 和 $1/8$，在每一爻中出现变爻的概率为 $1/2$。表1展示了在 4,194,304 次占卜结果中，变爻个数的分布情况：出现频率最高的为 1 个变爻（约 35.59\%），其次为 2 个变爻（约 29.65\%）和 0 个变爻（约 17.83\%）。变爻个数为 6 的情况极为罕见，仅占 0.025\%。该分布基本符合铜钱卦生成过程中的数学预期，说明随机数种子的选取在统计上具有良好的随机性与稳定性，不存在系统性偏差。

\begin{table}[]
\centering
\caption{变爻个数的分布}
\begin{tabular}{lcc}
\toprule
 & 出现频率 & 观测数 \\
\midrule
0个变爻 & 0.178275824 & 747743 \\
1个变爻 & 0.355926037 & 1492862 \\
2个变爻 & 0.296473265 & 1243499 \\
3个变爻 & 0.131716490 & 552459 \\
4个变爻 & 0.033016443 & 138481 \\
5个变爻 & 0.004341125 & 18208 \\
6个变爻 & 0.000250816 & 1052 \\
\bottomrule
\end{tabular}
\end{table}

起卦后，参考朱熹在《周易本义》中提倡的解卦法进行解卦。无变爻取本卦卦辞；一个变爻取本卦变爻爻辞；二个变爻取本卦中两变爻的爻辞，以上为主，另一为次；三个变爻视初爻是否变化决定以变卦卦辞为主还是以本卦卦辞为主；四个变爻取变卦中不变两爻的爻辞，以下为主，另一为次；五个变爻取变卦中不变爻的爻辞为主，另一变爻为次；六个变爻若为乾卦取“用九”、坤卦取“用六”，否则取变卦卦辞。

在 4096 种结果中，无变爻、三变爻和六变爻以卦辞及其组合解释，其余以爻辞及其组合解释。故首先依据民间对卦辞的划分进行量化，定义上上卦（乾1、坤2、比8等）为吉，中上卦（需5、师7等）、中中卦（泰11、否12等）与中下卦为平（蒙4、讼6等），下下卦为凶（屯3、小畜9等）。对于 384 爻，依据民间假托邵雍的分类方式划分吉/平/凶。对于“乾卦用九”和“坤卦用六”，则参考普遍研究结论，均视为吉利。

接下来，对于以两条爻辞或两条卦辞解释的结果，根据主次组合结果划分为 9 等（吉吉，吉平、吉凶、平吉、平平、平凶、凶吉、凶平、凶凶）。对于以一条爻辞或一条卦辞解释的结果，视其次要结果与主要结果相同，以相同规则划分结果等次。表2展示了九个等次结果的分布情况。由于本文的解释变量为等次，故采用有序离散模型进行回归。同时，在出生日期与长期运势的预测部分，本文仅使用主要结果做为吉凶的判断方式，将所有结果划分为 3 个等级，以进行稳健性讨论。

\begin{table}[]
\centering
\caption{占卜结果等次分布}
\begin{tabular}{lcc}
\toprule
 & 出现频率 & 观测数 \\
\midrule
得分9（吉吉） & 0.289300680 & 1213415 \\
得分8（吉平） & 0.044584274 & 187000 \\
得分7（吉凶） & 0.054158449 & 227157 \\
得分6（平吉） & 0.053043365 & 222480 \\
得分5（平平） & 0.204755306 & 858806 \\
得分4（平凶） & 0.042274952 & 177314 \\
得分3（凶吉） & 0.056972980 & 238962 \\
得分2（凶平） & 0.043227196 & 181308 \\
得分1（凶凶） & 0.211682796 & 887862 \\
\bottomrule
\end{tabular}
\end{table}

在考察幸运数字的偏好行为的部分，本文对 4,194,304 条随机数种子都取其个位数字，以是否为 0$\sim$9 为解释变量（$D$），依次纳入分析。在考察出生日期对于长期运势影响的部分，本文以该日期出生的儿童至入学截止日期的真实年龄（天数）为解释变量。区分1986年《中华人民共和国义务教育法》确立义务教育制度前后的样本，研究入学时真实年龄对占卜结果的影响。

为进一步考察结果的稳健性，本文在估计时增加了变爻个数的固定效应，这是由于在《周易》占筮的理论体系中，变爻的个数是衡量事物运动状态与变化剧烈程度的核心标尺。从哲学层面看，变爻象征着阴阳二气消长过程中的临界点，代表着事物发展到了必须转化或突破的关口。变爻个数本质上是对“时”与“势”的动态模拟——变爻越少，不确定性越小，趋势越明确；变爻越多，则变化越复杂，越需要从全局把握吉凶悔吝的转化规律。同时，在估计时使用聚类至变爻个数层次的聚类稳健标准误。

\section{回归结果}

\subsection{幸运数字的偏好与短期运势}

表3报告了幸运数字偏好对占卜结果的影响。结果显示，以数字“6”结尾的随机数种子显著降低了占卜结果的吉利程度（系数为-0.004，p<0.01），而其他数字的影响均不显著。这一发现支持了负性情绪状态下个体倾向于选择文化幸运数字的心理机制，尽管这种选择并未带来更有利的占卜结果。

\begin{table}[]
\centering
\caption{随机数种子的个位数字与占卜结果}
\begin{tabular}{lccccc}
\toprule
个位 & (0) & (1) & (2) & (3) & (4) \\
\midrule
D & $0.003^{*}$ & $-0.000$ & $-0.002$ & $0.002$ & $-0.001$ \\
  & $[0.00]$ & $[0.00]$ & $[0.00]$ & $[0.00]$ & $[0.00]$ \\
N & 4194304 & 4194304 & 4194304 & 4194304 & 4194304 \\
Fixed & Yes & Yes & Yes & Yes & Yes \\
\bottomrule
\end{tabular}

\vspace{1em}

\begin{tabular}{lccccc}
\toprule
个位 & (5) & (6) & (7) & (8) & (9) \\
\midrule
D & $-0.000$ & $-0.004^{***}$ & $0.001$ & $0.000$ & $0.001$ \\
  & $[0.00]$ & $[0.00]$ & $[0.00]$ & $[0.00]$ & $[0.00]$ \\
N & 4194304 & 4194304 & 4194304 & 4194304 & 4194304 \\
Fixed & Yes & Yes & Yes & Yes & Yes \\
\bottomrule
\end{tabular}
\end{table}

\subsection{出生日期与长期运势}

表4聚焦于出生日期对占卜结果的影响。在1986年《中华人民共和国义务教育法》实施后，入学真实年龄每增加一天，占卜结果的吉利程度显著提高，而在1986年前，这一关系为负向。这表明，义务教育制度实施后，出生日期对个体长期发展的影响更加显著，进而体现在占卜结果的吉利程度上。

\begin{table}[H]
\centering
\caption{出生日期与长期运势}
\begin{tabular}{lcccc}
\toprule
 & \multicolumn{2}{c}{1986年之后} & \multicolumn{2}{c}{1986年之前} \\
\cmidrule(lr){2-3} \cmidrule(lr){4-5}
划分方式 & 9等划分 & 3等划分 & 9等划分 & 3等划分 \\
\midrule
真实年龄 & $0.000^{***}$ & $0.000^{***}$ & $-0.000^{*}$ & $-0.000^{**}$ \\
         & $[0.00]$      & $[0.00]$      & $[0.00]$     & $[0.00]$ \\
N        & 16924         & 16924         & 19601        & 19601 \\
Fixed    & Yes           & Yes           & Yes          & Yes \\
\bottomrule
\end{tabular}
\end{table}

\subsection{稳健性分析}

出于对辩证唯物主义与历史唯物主义的坚定认同，笔者认为本文结果在理论上并不稳健，只是统计上特别显著，故不宜进行进一步讨论。

\section{结论}

本文通过对chgdiv6命令生成的超过400万条占卜结果进行实证分析，探讨了“赛博算卦为何很准”这一现象背后的统计规律。研究发现，求卦者在选择随机数种子时无意中融入了自身的心理状态与生活背景，进而影响了占卜结果的分布。具体而言，负性情绪下对幸运数字“6”的偏好显著降低了占卜结果的吉利程度；而在义务教育制度实施后，出生日期所决定的入学真实年龄与占卜结果的吉利程度呈正相关，反映了出生时间对个体长期发展的潜在影响。

尽管从设计角度看，chgdiv6命令仅是对随机数生成过程的文化包装，但其结果之所以被认为“准”，部分原因在于求卦者自身特征与心理状态通过随机数种子的选择无意识地嵌入了占卜过程。因此，占卜的“准确性”并非源于超自然力量，而是源于人类行为的可统计规律性。

% --- 输出参考文献 ---
\erjref
\printbibliography[heading=none]

\end{document}